% Extracto íntegro por residencia del propietario científico definitivo de masas.
\section{Null sector and positivity}

Nullity is not obtained by taking a limit of the positive character. The domain
is first separated into
\[
 \mathcal H_{\rm phys}=\mathcal H_0\oplus\mathcal H_+,
\]
with
\[
 \mathcal M|_{\mathcal H_0}=0,
 \qquad
 \mathcal M|_{\mathcal H_+}>0.
\]
The photon belongs to the electromagnetic null chart and the gluons to the adjoint
null chart of color. Their different representations, polarizations, and
gauge restrictions remain visible even though both masses are null.

Positivity excludes an elementary species with \(m^2<0\) within the
mass codomain. A negative eigenvalue of the Hessian of a background does not define a
tachyon as a particle: it diagnoses the instability of the background and requires changing
the solution around which linearization is performed.

\section{Color, quark flavors, and CKM mixing}

The chromatic orbit is realized in the color qutrit. The Weyl operators
generate \(\mathfrak{sl}_3(\mathbb C)\); the traceless anti-Hermitian selection
yields \(\mathfrak{su}(3)\), and the finite connection transports the three
chromatic charts. The flavor mass must be invariant under this action.

Let \(U_{\rm col}:SU(3)\to\mathcal U(V_{\rm col})\). The chromatic reduction
is the Reynolds expectation
\[
 \mathbb E_{\rm col}(T)
 =\int_{SU(3)}U_{\rm col}(g)T U_{\rm col}(g)^*\,dg.
\]
In the fundamental irreducible representation,
\[
 \mathbb E_{\rm col}(\mathcal M_q)
 =m_q(\mu,\mathscr R)I_3.
\]
The published mass is therefore independent of red, green and blue. The
coordinates \(\mu\) and \(\mathscr R\) identify the scale and renormalization
scheme; they are not hidden ambiguities and must always accompany the
figure.

The dependence on scale is a section of a fibration over
\(t=\log\mu\). If \(\gamma_m(g(t))\) is the endomorphism induced by the
chromatic connection, transport is defined by
\[
 \nabla_{\rm RG}m_q
 =\left(\frac{d}{dt}+\gamma_m(g(t))\right)m_q=0,
\]
\[
 m_q(t_2)
 =\mathcal P\exp\!\left(-\int_{t_1}^{t_2}
 \gamma_m(g(t))\,dt\right)m_q(t_1).
\]
Two values belonging to different scales or schemes are not comparable
numbers until they are transported to the same fiber. This equation closes the
mathematical type of the running mass: it is not an isolated absolute scalar, but
a covariantly transported section.

For mesons and baryons, the singlet projectors are
\[
 P_{q\bar q}^{(1)}
 =\frac1{3}|\Omega\rangle\langle\Omega|,
 \qquad
 |\Omega\rangle=\sum_{a=1}^{3}|a\bar a\rangle,
\]
and
\[
 P_{qqq}^{(1)}
 =|\varepsilon\rangle\langle\varepsilon|,
 \qquad
 |\varepsilon\rangle
 =\frac1{\sqrt6}\sum_{a,b,c}\varepsilon_{abc}|abc\rangle.
\]
Both eliminate the observable chromatic label without erasing the internal orbit
that intervened in the construction.

The CKM mixing belongs to a distinct map. If \(B_u\) and \(B_d\) are the
mass operators in the up-type and down-type bases, the internal unitary matrix
satisfies
\[
 B_d^{(u)}=V_{\rm CKM} B_d V_{\rm CKM}^*.
\]
The angles derived from \(A\), \(C^*\) and \(\Delta_4\), the phase
\(\delta_{\rm CP}\), the Jarlskog invariant and the determinant identity
of the commutator determine mixing without creating the mass eigenvalues. Thus
spectrum, color, generation and change of basis remain separate.

\section{Neutrinos, PMNS Mixing, and Majorana Criterion}

Neutrinos are electrically neutral fermionic sections. Their
fermionicity descends from the spinorial holonomy, the exterior completion, and
the canonical anticommutation relations; neutrality does not alter their
statistics. The internal support is
\[
 \mathcal H_\nu=
 \bigoplus_{j=1}^{4}\mathcal H_{\nu,G_j},
 \qquad
 \dim\mathcal H_\nu=2+4+6+8=20.
\]
The four structural depths are neither flavours nor eigenvalues. The
three-dimensional physical subspace is obtained through the neutral projector
\(P_\nu\), and the mass operator is
\[
 M_\nu=P_\nu\mathcal M_\nu P_\nu.
\]
Its three ordered eigenvalues are the masses; the absolute scale comes from the
same mass clock that fixes the remaining branches, and not from the mixing angles.

Let \(M_\ell=P_\ell\mathcal M_\ell P_\ell\) be the charged operator. Denote by
\(F_\ell\) and \(F_\nu\) complete orthonormal spectral frames, ordered by
the spectral projectors of \(M_\ell\) and \(M_\nu\). The mixing matrix is
\[
 U_{\rm PMNS}=F_\ell^*F_\nu.
\]

\begin{theorem}[Spectral closure of the leptonic mixture]
If the two frames are complete, \(U_{\rm PMNS}\) is unitary. If the
spectra are simple, it is determined up to diagonal phase readjustments and
permutations fixed by the spectral order. If degeneracy exists, the
exact result is a double class
\[
 \prod_\lambda U(m_{\ell,\lambda})
 \backslash U(3) /
 \prod_\mu U(m_{\nu,\mu}),
\]
and not a numerically artificial matrix selected.
\end{theorem}

\begin{proof}
Orthonormality gives
\(U_{\rm PMNS}^*U_{\rm PMNS}=F_\nu^*F_\ell F_\ell^*F_\nu=I_3\).
Changes of basis within each degenerate subspace act on the
left or on the right, producing exactly the indicated
double class.
\end{proof}

The complete record determines the elements of \(M_\nu\) by means of carry,
sheet, orientation, vacancy, and holonomy. Equivalently, the
Hermitian kernel may be used
\[
 K_{\rm reg}(c,d)
 =\sum_{r\in\{\rm car,hoj,ori,vac,hol\}}
 \frac{\langle\Xi_r(c),\Xi_r(d)\rangle}
      {\sum_x\|\Xi_r(x)\|^2},
\]
with omission of every summand of null norm. Its Gram matrix is positive and
constructs the operator before diagonalization. The rank-one phase
\[
 I+(e^{i\pi/6}-1)|b_K\rangle\langle b_K|
\]
remains as a negative control: its formal unitarity is not sufficient and its angle
\(\theta_{13}\) falsifies its identification with full-rank physical
mixing.

Let \(G_0\) be the three-dimensional matrix obtained by aggregating the preceding kernel
over the three neutrino sectors. The finite condition
\[
 \det G_0\ne0
\]
is the exact certificate that the observables included in the record
separate three directions. In that case, the polar factor
\[
 U_{G_0}=G_0(G_0^*G_0)^{-1/2}
\]
is unitary and has rank \(3\). If \(\det G_0=0\), the declared kernel does not
determine a three-dimensional transport: an independent internal observable
must be added and the Gram matrix recomputed. Invertibility may not be
imposed by adding a term constructed from the mixing matrix one seeks to obtain.
The physical matrix is fixed by the spectral frames of \(M_\ell\) and
\(M_\nu\); the Gram criterion certifies sufficiency of the register, but does not
by itself select those frames.

Let \(\mathcal C_\nu\) be the antiunitary self-conjugation of the neutral sector. The
Majorana criterion is
\[
 \mathcal C_\nu^2=I,
 \qquad
 \mathcal C_\nu M_\nu=M_\nu\mathcal C_\nu,
 \qquad
 \mathcal C_\nu H_\nu=H_\nu\mathcal C_\nu,
\]
together with the existence of a fixed real basis in the realification of the
physical subspace. When these conditions are satisfied, each mode can be
represented by \(\psi=\mathcal C_\nu\psi\). The Dirac branch requires, by
contrast, a conserved charge \(Q\) such that
\[
 \mathcal C_\nu Q\mathcal C_\nu^{-1}=-Q,
 \qquad Q\psi\ne0,
\]
and retains \(\psi\) and \(\mathcal C_\nu\psi\) as distinct modes. If neither
set of conditions is established, the type remains undecided.
Neutrality, the word involution or the Ising sector, taken
separately, do not decide this alternative.

An additional structural depth does not by itself constitute a sterile
neutrino. A candidate sterile sector is defined by a nonzero projector
\(P_s\) satisfying
\[
 P_s P_{\rm act}=0,
 \qquad P_s J_{\text{weak}}=0,
\]
with dynamics defined on \(P_s\mathcal H_\nu\). If, in addition,
\([P_s,M_\nu]=0\), the sector is reducing and does not mix with the active sector. If
\(P_sM_\nu P_{\rm act}\ne0\), active--sterile mixing exists and necessarily
\([P_s,M_\nu]\ne0\). Thus, existence of the sector and existence of
mixing are distinct properties. If the weakly neutral projector does not exist,
the additional coarse channel is internal depth of the support and not a fourth
species.

\section{Massive bosons and the scalar sector}

The sectors \(W\), \(Z\), and the scalar sector have their own
representation, route, and signature descriptors. Their realization does not require forcing them into an
elementary rank-one fiber. The positive fiber operator is constructed first
and then, on a finite-dimensional subspace \(P_X\mathcal H_X\), the
effective channel operator
\[
 H_{X,\rm eff}(z)
 =P_X H_X P_X
 +P_X H_X Q_X(z-Q_X H_X Q_X)^{-1}Q_X H_X P_X.
\]
If the reduced resolvent admits meromorphic continuation through the cut of the
continuum, the poles are the isolated zeros, counted with multiplicity, of
\[
 \det\bigl(z-H_{X,\rm eff}(z)\bigr)=0,
 \qquad
 z_X=M_X c_{\rm int}^2-\frac{i}{2}\Gamma_X.
\]
An open channel does not by itself guarantee the existence of such zeros. When
a zero with the prescribed physical residue exists, its real part publishes pole
mass and its imaginary part the width. For the
scalar sector, identification of an excitation additionally requires the scalar
projector and its gauge and fermionic couplings. Two independent poles
determine two states; a second label without projector or pole does not create a
particle.

\section{Composite states and binding law}

For constituents \(X_1,\ldots,X_n\), the input is the composite record
\[
 \mathfrak C_{12}^{(n)}:
 (\Gamma_1,\ldots,\Gamma_n;\mathcal T,\partial)
 \longmapsto\mathcal L_{\rm comp},
\]
not the sum of isolated masses. The binary composition
\[
 L_1\star L_2=L_1+L_2+\mathfrak b(L_1,L_2)
\]
is associative exactly when
\[
 \mathfrak b(L_1,L_2)+\mathfrak b(L_1\star L_2,L_3)
 =\mathfrak b(L_2,L_3)+\mathfrak b(L_1,L_2\star L_3).
\]
The signature, the refinement, and the composite operator are computed after the linking:
\[
 \mathcal L_{\rm comp}\longmapsto\sigma_{\rm comp}
 \longmapsto\eta_{\rm comp}\longmapsto\mathcal M_{\rm comp}.
\]
The physical projector combines color singlet, spin and parity, isospin, charge,
exchange symmetry, and channel. The stable mass is an eigenvalue of the
compression; when the hypotheses of meromorphic continuation hold, the mass
and width of a resonance follow from a Feshbach pole. Thus,
mesons, baryons, exotic states, nuclei, atoms, and
molecules possess the same causal scheme without being reduced to a sum of masses.

Two constituent trees represent the same species if there exists a
unitary that simultaneously intertwines their operators, channel projectors and
poles. In the absence of such a unitary they are retained as distinct realizations,
even if they share constituents and partial quantum numbers.

\section{Resonances, widths, and survival}

Let \(S_X\) be the projector onto the surviving subspace and
\(T_X=S_X U S_X\) the compressed one-step operator. If it has a simple dominant
eigenvalue
\[
 \lambda_X=r_X e^{-i\theta_X},\qquad 0<r_X<1,
\]
fix the branch \(\theta_X\in I_\theta\), where \(I_\theta\) is the fundamental
quasienergy interval determined by the temporal chart and its orientation.
Then
\[
 M_X=\frac{\hbar_{\rm int}}{t_0c_{\rm int}^2}\theta_X,
 \qquad
 \Gamma_X=-\frac{2\hbar_{\rm int}}{t_0}\log r_X,
 \qquad
 \tau_X=\frac{\hbar_{\rm int}}{\Gamma_X}.
\]
Without the choice of \(I_\theta\), the phase determines only \(\theta_X\) modulo
\(2\pi\), and the quasienergy mass is not a well-defined scalar.
The width is not a sixth coordinate of the mass signature: it arises from the modulus
of the open mode. Equality of mass and width between particle and
antiparticle also requires transporting the channels and the causal prescription, not merely
the parity of the signature.

\section{Multiparticle completion and statistics}

Let \(\mathsf{Hol}_{\pm}^{\leq1}\) be the monoidal category of fibers of a
particle, with contractive refinement morphisms, in particular unitary ones,
and holonomy parity
\(\epsilon\in\{+1,-1\}\). The Fock completion is defined by the functor
\[
 \mathfrak F(H,\epsilon)=
 \begin{cases}
 \displaystyle\bigoplus_{n\ge0}\operatorname{Sym}^n(H),
     &\epsilon=+1,\\[2mm]
 \displaystyle\bigoplus_{n\ge0}\bigwedge^n\!H,
     &\epsilon=-1,
 \end{cases}
\]
and, for a contractive intertwiner \(T:H\to H'\),
\[
 \mathfrak F(T)=
 \bigoplus_{n\ge0}\operatorname{Sym}^n(T)
 \quad\text{or}\quad
 \bigoplus_{n\ge0}\bigwedge^n\!T.
\]

\begin{theorem}[Internal transport of the statistic]
The foregoing completion is a functor of bounded operators on \(\mathsf{Hol}_{\pm}^{\leq1}\); it is compatible with refinement and is graded monoidal. In the \(\epsilon=-1\) branch, the creation and annihilation operators satisfy CAR and the occupation operators satisfy \(N_s^2=N_s\); in the \(\epsilon=+1\) branch, they satisfy CCR. Therefore, Pauli exclusion and fermionic and bosonic statistics descend from holonomy parity once the representation of a particle has been fixed.
\end{theorem}

\begin{proof}
Symmetric and exterior powers preserve compositions and identities; the contractivity of \(T\) ensures that their direct sum defines a bounded operator. For noncontractive morphisms, second quantization requires an explicit dense domain and does not belong to this category of bounded operators. The sign rule of the exterior product produces CAR and idempotence of occupation; the symmetric quotient produces CCR. Second quantization of the refinement morphisms intertwines the corresponding operators.
\end{proof}

This proof closes the internal algebraic transport. Identification
with the relativistic spin--statistics theorem additionally requires the
Lorentzian bridge and locality declared in Dimensional Mechanics; the two statements
must not be confused.

For a compressed Hamiltonian \(H_{X,V}\) over a finite volume or refinement
\(V\), the thermal state is completely defined by
\[
 Z_{X,V}(\beta)=\operatorname{Tr}e^{-\beta H_{X,V}},
 \qquad
 \rho_{X,V,\beta}=Z_{X,V}(\beta)^{-1}e^{-\beta H_{X,V}}.
\]
The thermodynamic limit exists when the pressure
\[
 p_X(\beta)=\lim_{V\nearrow\infty}
 \frac1{\beta |V|}\log Z_{X,V}(\beta)
\]
exists and is independent of the cofinal sequence. Bosonic condensation is decided by saturation of the density of excited states; it is not proclaimed from the mere existence of a symmetric completion. The problem is thus expressed through a verifiable mathematical condition for the Hamiltonian and the limit, rather than through a nominal attribution.

\section{Topological sectors and virtual states}

In a topological sector the primary codomain is a braided category, not
a line of masses. For simple objects \(a,b,c\), the physical data include
\[
 N_{ab}^{c},\qquad F^{abc}_{d},\qquad R^{ab}_{c},
\]
with pentagon and hexagon equations. A scalar mass exists only
when a Hamiltonian realization provides a spectral projector and a
gap. The Fibonacci, Ising/Majorana, and \(\mathbb Z_6\) sectors are thereby
closed as categorical codomains; they are not assigned a universal figure
by analogy with an elementary particle.

In the sector \(\mathbb Z_6\), fusion and braiding are
\[
 a\otimes b=a+b\pmod 6,
 \qquad
 R_{ab}=\zeta_6^{ab},
 \qquad \zeta_6=e^{2\pi i/6}.
\]
For Fibonacci,
\[
 \tau\otimes\tau=\mathbf 1\oplus\tau,
 \qquad
 F_\tau^{\tau\tau\tau}=
 \begin{pmatrix}
 \varphi^{-1}&\varphi^{-1/2}\\
 \varphi^{-1/2}&-\varphi^{-1}
 \end{pmatrix},
\]
\[
 R_{\mathbf1}^{\tau\tau}=e^{-4\pi i/5},
 \qquad
 R_\tau^{\tau\tau}=e^{3\pi i/5}.
\]
The Ising sector has objects \(\mathbf1,\psi,\sigma\) and rules
\[
 \psi\otimes\psi=\mathbf1,
 \qquad
 \psi\otimes\sigma=\sigma,
 \qquad
 \sigma\otimes\sigma=\mathbf1\oplus\psi,
\]
\[
 F_\sigma^{\sigma\sigma\sigma}
 =\frac1{\sqrt2}
 \begin{pmatrix}1&1\\1&-1\end{pmatrix},
 \qquad
 R_{\mathbf1}^{\sigma\sigma}=e^{-\pi i/8},
 \qquad
 R_\psi^{\sigma\sigma}=e^{3\pi i/8}.
\]
These matrices satisfy the coherence relations in the indicated
conventions. The duplication
\(\mathcal C\boxtimes\mathcal C^{\rm rev}\) preserves both orientations and
avoids introducing a chirality by external choice.

A virtual state is a section of finite persistence in an amplitude or a
resolvent. It does not belong to the asymptotic spectrum and does not require a real pole
mass. Its parameters are virtuality, the channel and the duration of
persistence. This definition avoids turning an internal computational line into
a new species.

\Needspace{25\baselineskip}
\section{Gravity, torsion, and spin-2 mode}

Gravitation enters Dimensional Mechanics by means of connection, curvature, and
torsion. It does not require an elementary particle of spin \(2\) as a datum of
departure. Let \(\mathcal S_{\rm grav}\) be the geometric action and
\(\mathcal H_{\rm TT}\) the transverse traceless subspace of its Hessian
in a solution. The linearized observable is
\[
 \mathcal K_{\rm TT}
 =P_{\rm TT}\,\delta^2\mathcal S_{\rm grav}\,P_{\rm TT}.
\]
A collective spin \(2\) mode exists as an asymptotic particle only
if the resolvent of \(\mathcal K_{\rm TT}\) has a physical pole with positive
residue and satisfies the constraints. If no such pole exists, the gravitational
interaction remains wholly in the connection and torsion. Thus, the
graviton is not a necessary component of the mass law, and its possible appearance
is decided by a spectral criterion, not by a nominal box.
